\chapter{Vitiligo}
\index{Vitiligo}
\label{sec:Vitiligo}

\section{Overview}
Vitiligo is an acquired autoimmune disorder characterized by selective destruction of melanocytes, resulting in well-demarcated, milky-white macules and patches, affecting 0.5--2\% of the global population.
\section{Causes}
This condition results from disruption of normal vitiligo function.

\begin{itemize}[noitemsep]
  \item This condition happens because of autoimmune-mediated melanocyte destruction. Loss of pigment may occur in any area but commonly affects the face, hands, elbows, knees, and genitalia. Vitiligo is linked to other autoimmune diseases (thyroid disease, type 1 diabetes, Addison's disease, pernicious anemia)
  \item The Koebner phenomenon (new lesions at sites of trauma) may occur.
\end{itemize}
\section{Risk Factors}


\begin{itemize}[noitemsep]
  \item Genetic predisposition and family history
  \item Advancing age
  \item Lifestyle factors (diet, exercise, smoking, alcohol)
  \item Environmental exposures
  \item Underlying medical conditions
  \item Medication use
\end{itemize}
\section{Signs and Symptoms}

\subsection{Common symptoms}
\begin{itemize}[noitemsep]
  \item Pain or discomfort in the affected area
  \end{itemize}

\subsection{Emergency warning signs}
\begin{itemize}[noitemsep]
  \item Severe or worsening symptoms of vitiligo require immediate medical evaluation
\end{itemize}
\section{Progression and Stages}
It may be segmental (unilateral, dermatomal, often with early onset and stable course) or non-segmental (generalized, bilateral, with variable progression). The condition may progress through different stages of severity over time. Early stages often involve mild or subtle symptoms that may be overlooked. Moderate stages involve more noticeable symptoms affecting daily function. Advanced stages may involve significant impairment and complications.
\section{Complications}
\begin{itemize}[noitemsep]
  \item Disease progression leading to worsening symptoms
  \item Secondary infections or injuries
  \item Side effects of treatment
  \item Reduced quality of life
  \item Emotional and psychological impact
\end{itemize}
\section{Diagnosis}
Doctors usually diagnose this based on your symptoms and a physical exam, supported by Wood's lamp examination. Comprehensive medical history and physical examination Laboratory tests to assess organ function and rule out other conditions Imaging studies to visualize affected structures Specialized testing depending on the affected system Consultation with relevant specialists Monitoring of disease progression over time

\begin{itemize}[noitemsep]
  \item Comprehensive medical history and physical examination
  \item Laboratory tests to assess organ function
  \item Imaging studies to visualize affected structures
  \item Specialized testing depending on the affected system
  \item Consultation with relevant specialists
\end{itemize}
\section{Prevention}
\begin{itemize}[noitemsep]
  \item Maintain a healthy lifestyle with balanced diet and exercise
  \item Avoid known risk factors
  \item Attend regular medical check-ups for early detection
  \item Follow recommended screening guidelines
\end{itemize}
\section{Treatment Options}
\begin{itemize}[noitemsep]
  \item Treatment options include topical corticosteroids, topical calcineurin inhibitors (tacrolimus, especially for facial lesions), narrowband UVB phototherapy, and for refractory localized disease, surgical melanocyte transplantation
  \item Ruxolitinib cream (JAK1/2 inhibitor) has been approved for non-segmental vitiligo. Sun protection is essential.
  \item Medications to manage symptoms and address underlying mechanisms
  \item Lifestyle modifications and self-care measures
  \item Physical therapy and rehabilitation
  \item Surgical intervention when indicated
  \item Regular monitoring and follow-up care
\end{itemize}
\section{Living with It}
\begin{itemize}[noitemsep]
  \item Follow your treatment plan as prescribed
  \item Maintain regular follow-up with your healthcare provider
  \item Make necessary lifestyle adjustments
  \item Seek support from family, friends, or support groups
  \item Stay informed about your condition
\end{itemize}
\section{Outlook}
\begin{itemize}[noitemsep]
  \item The outlook varies from person to person
  \item Repigmentation is achievable but recurrence may occur.
\end{itemize}
